\section{问题二：仅知当天天气的反馈决策}

\subsection{信息约束与概率边界}

设$H_t=(S_0,z_1,a_1,\ldots,S_{t-1},z_t)$为行动前完整历史，$\mathcal F_t=\sigma(H_t)$。在线策略$\pi_t(H_t)$必须满足\cref{eq:nonanticipativity}；代码层面只向策略对象传入历史记录和当天气象，未来数组封装在外部环境中，因而信息约束不是文字承诺，而是接口约束。天气概率并非题面事实，本文作如下条件设定：第三关无沙暴，以高温概率$p_h\in\{0.2,0.5,0.8\}$的独立分布做精确敏感性；第四关分别使用
\[
(p_s,p_h,p_b)=(0.50,0.45,0.05),\ (0.45,0.45,0.10),\ (0.35,0.50,0.15)
\]
代表低沙暴、基准和不利天气。

\subsection{第三关有限时域精确 MDP}

第三关没有村庄，现金只受初始采购、挖矿和终点回收影响。给定初购后，历史对未来的作用可由充分状态$s=(l,w,f)$概括。为准确体现“先观测天气再行动”，分别定义天气揭示前价值$V_t(s)$和揭示后价值$Q_t(s,z)$：
\begin{align}
Q_t(s,z)&=\max_{a\in\mathcal A(s,z)}
\left\{R(s,a,z)+V_{t+1}\bigl(F(s,a,z)\bigr)\right\},
\label{eq:post-weather-bellman}\\
V_t(s)&=\sum_{z\in\mathcal Z}p_zQ_t(s,z).
\label{eq:pre-weather-bellman}
\end{align}
其中即时收益$R=B$当且仅当挖矿，否则为0；$F$按\cref{eq:resource-transition}扣减资源并更新位置。边界条件为
\begin{equation}
V_t(v_T,w,f)=\tfrac12p_ww+\tfrac12p_ff,\qquad
V_{T+1}(l,w,f)=-\infty\quad(l\ne v_T).
\label{eq:mdp-boundary}
\end{equation}
现金在途中不影响行动可行性，故可从状态中剔除，并在初购层补回。完整优化为
\begin{equation}
\max_{w_0,f_0\in\mathbb Z_+}
\left\{M_0-p_ww_0-p_ff_0+V_1(v_S,w_0,f_0)\right\}
\quad\text{s.t. }3w_0+2f_0\le C.
\label{eq:mdp-initial-load}
\end{equation}
若某个$p_z>0$而$Q_t(s,z)=-\infty$，则令$V_t(s)=-\infty$。这一“支撑集稳健可行”递推排除了以小概率必败换取较高条件均值的方案。

有限时域和整数资源使状态空间有限。采用记忆化递归后，每个$(t,l,w,f)$只计算一次；同一状态、同一资源消耗下仅保留收益最大的行动。最坏复杂度为$O(T|V|WF|\mathcal Z|\mathcal A|)$，其中$W,F$分别是容量给出的水、食物整数上界。

程序枚举全部整数初始采购和可达状态。每档高温概率评估6425464个状态，并对全部 $2^{10}=1024$ 条天气序列逐条回放，Bellman期望与枚举加权均值误差小于 $10^{-6}$。

\begin{table}[H]
\centering
\caption{第三关精确 MDP 对高温概率的敏感性}
\label{tab:level3}
\begin{tabular}{ccrrrr}
\toprule
高温概率 & 初始水/食物 & 期望终值 & 5\%分位 & 最坏值 & 失败序列\\
\midrule
0.2 & 54/54 & 9382 & 9270 & 9190 & 0/1024\\
0.5 & 54/54 & 9310 & 9190 & 9190 & 0/1024\\
0.8 & 54/54 & 9238 & 9190 & 9190 & 0/1024\\
\bottomrule
\end{tabular}
\end{table}

三档概率下策略均为$1\to4\to6\to13$连续移动三天。全高温路径的最低资源由
\begin{equation}
w_0^{\min}=3\times2c^w_{\text{高温}}=54,\qquad
f_0^{\min}=3\times2c^f_{\text{高温}}=54
\label{eq:level3-minload}
\end{equation}
给出，恰好等于MDP解。绕矿需要五个移动日，矿山基础收益仅200元，额外稳健装载和绕行成本超过挖矿收益。因此“54/54直达”既有Bellman最优性，也有由最坏支撑资源下界给出的结构解释。

\begin{figure}[H]
  \centering
  \includegraphics[width=0.84\textwidth]{level3_mdp.pdf}
  \caption{第三关高温概率、期望终值与支撑内风险}
  \label{fig:level3}
\end{figure}

\subsection{第四关鲁棒阈值策略}

第四关状态空间远大于第三关，且含沙暴、矿山和较长时域，直接精确DP在多档参数搜索中代价过高。我们从“挖矿后仍必须保证返程”的充分条件出发构造可解释阈值策略。设$d_t=d(l_t,v_T)$为当前位置到终点最短边数，$N_b(t)$为截至当天已观察到的沙暴数，$B$为设计的总沙暴预算，剩余预算为
\begin{equation}
\widehat B_t=(B-N_b(t))_+.
\label{eq:remaining-storm-budget}
\end{equation}
若以后最多再遇$\widehat B_t$个沙暴，则$d_t$个移动日按高温上界需$2d_tc_h^r$，沙暴停留需$\widehat B_tc_b^r$。因此第$t$日挖矿后保持返程可行的充分条件为
\begin{align}
T-t &\ge d_t+\widehat B_t,\label{eq:threshold-time}\\
w_{t-1}-3c^w_{z_t} &\ge 2d_tc^w_{\text{高温}}+\widehat B_t c^w_{\text{沙暴}},
\label{eq:threshold-water}\\
f_{t-1}-3c^f_{z_t} &\ge 2d_tc^f_{\text{高温}}+\widehat B_t c^f_{\text{沙暴}}.
\label{eq:threshold-food}
\end{align}
式中右端第一项的系数2来自移动消耗倍数，第二项的系数1来自沙暴停留。三式任一不满足时，非沙暴日立即沿最短路向终点移动，沙暴日停留。阈值不是对所有天气路径的绝对保证，而是把风险偏好压缩到一个可搜索、可解释的整数$B$。

\subsection{567候选选模与独立测试}

枚举$B=3,\ldots,9$，初始水食180--280箱、步长10，并过滤负重与现金不合法组合，得到567个候选$\Theta$。对每个$\theta\in\Theta$，在基准和不利分布各使用1000条相同的配对天气序列，减少候选比较中的抽样噪声。训练选择规则在模拟前固定为
\begin{equation}
\widehat\theta=\arg\max_{\theta\in\Theta}\widehat\mu_{\mathrm{adv}}(\theta)
\quad\text{s.t.}\quad
\widehat p_{\mathrm{base}}(\theta)\le0.002,
\quad
\widehat p_{\mathrm{adv}}(\theta)\le0.005,
\label{eq:level4-selection}
\end{equation}
其中$\widehat\mu$按\cref{eq:penalized-objective}把失败样本记0。入选水240箱、食物240箱、$B=5$。\cref{fig:level4-frontier}显示，入选点不是训练均值最大的点，而是风险预算内的收益折中。

\begin{figure}[H]
  \centering
  \includegraphics[width=0.98\textwidth]{level4_frontier.pdf}
  \caption{第四关候选策略在基准和不利天气下的训练收益—风险前沿}
  \label{fig:level4-frontier}
\end{figure}

最终测试改用从未参与选模的种子20260821，每种天气生成20000条新序列。设$I_j$为第$j$条序列的失败指示，$Y_j$为失败记0后的终值，则报告
\begin{equation}
\widehat p=\frac1N\sum_{j=1}^NI_j,\qquad
\widehat\mu=\frac1N\sum_{j=1}^NY_j.
\label{eq:oos-estimators}
\end{equation}
失败率使用Wilson区间，均值使用样本标准误差区间；完整公式和图形检验见附录。

\begin{table}[H]
\centering
\caption{第四关入选策略的独立样本外结果}
\label{tab:level4-oos}
\begin{tabular}{lrrrr}
\toprule
天气模型 & 失败率 & 95\% Wilson区间 & 总体均值/元 & 总体5\%分位/元\\
\midrule
低沙暴 & 0 & [0, 0.0192\%] & 10358.8 & 8880.0\\
基准 & 0.055\% & [0.0307\%, 0.0985\%] & \textbf{10162.9} & 8805.0\\
不利 & 0.385\% & [0.3082\%, 0.4809\%] & 9635.3 & 8627.5\\
\bottomrule
\end{tabular}
\end{table}

基准天气下，入选B5比原B6策略总体均值高418.3元；不利天气下高391.9元，并保持失败率上置信界0.481\%低于0.5\%预算。B4虽有更高成功者均值，但基准与不利失败率分别升至0.145\%和1.005\%，不符合不利分布预算。

\begin{resultbox}{问题二结论}
第三关选择54/54箱三日直达，属于全支撑精确结论\evidence{E2}。第四关选择240/240箱、B5挖矿阈值策略，属于明确分布和风险预算下的样本外推荐\evidence{E4}；二者的“最优”含义不同。
\end{resultbox}
